% TEMPLATE-MILC-v3.tex - plantilla maestra UNICA de las guias MILC v3.
%
% La usan las tres familias de guias, cada una con su driver:
%   · Grados 6-11  -> scripts/build-guias-g11.py
%   · Bebras       -> scripts/build-guias-bebras.py
%   · Semillero    -> scripts/build-guias-semillero.py
%
% Antes habia tres copias casi identicas (~400 lineas iguales, 6 distintas):
% cada mejora habia que aplicarla tres veces, y en la practica no se hacia
% (el motor de recursos vivia solo en dos de ellas). Lo que varia entre
% familias esta parametrizado en 6 placeholders que cada driver llena
% (se nombran aqui SIN su sintaxis real: el builder reemplaza por texto plano,
% tambien dentro de los comentarios, y un valor multilinea rompe el preambulo):
%
%   HEADER_IZQ        encabezado izquierdo de pagina
%   HEADER_DER        encabezado derecho de pagina
%   PORTADA_ETIQUETA  rotulo sobre el numero grande (GRADO/MOMENTO/MODULO)
%   PORTADA_SERIE     linea de serie de la portada
%   PORTADA_ID        identificador de la guia en la portada
%   PORTADA_CIRCULO   el circulito de grado (°), o vacio si no aplica
%
% El resto de claves las documenta content/guias/_SCHEMA.md. El script valida
% que no queden placeholders sin reemplazar antes de escribir el archivo final.

\documentclass[11pt,letterpaper]{article}
\usepackage[margin=1.75cm]{geometry}
\usepackage[table]{xcolor}
\usepackage{fontspec}
\usepackage[spanish]{babel}
\usepackage{tikz}
% Librerías TikZ para los diagramas de las guías (bloque `recursos:`):
%  · babel  → babel en español vuelve activos caracteres como «>», y TikZ los usa
%             en sus opciones (>=stealth, ->). Sin esta librería, cualquier
%             diagrama con flechas revienta con «\language@active@arg> has an
%             extra }». Debe ir ANTES de que se dibuje nada.
%  · positioning        → sintaxis «below left=2pt and 3pt».
%  · decorations.pathreplacing → llaves ({brace}) para señalar rangos.
\usetikzlibrary{babel,positioning,decorations.pathreplacing}
\usepackage{graphicx}
% Circuitos: el trabajo de electrónica se hace hoy con micro:bit y sus sensores
% integrados (MakeCode, sin cableado), así que no se necesitan esquemáticos.
% Si algún día entran montajes con protoboard, basta descomentar esta línea:
% los diagramas .tex/.tikz declarados en `recursos:` ya se incrustan solos.
% \usepackage{circuitikz}
\usepackage[most]{tcolorbox}
\usepackage{tabularx}
\usepackage{array}
\usepackage{fancyhdr}
\usepackage{titlesec}
\usepackage{ragged2e}
\usepackage{enumitem}
\usepackage{microtype}

% Helvetica Neue no trae la flecha U+2192, asi que cada «→» escrito literalmente
% en el YAML se compilaba a NADA, en silencio (462 flechas en 61 guias: rutas del
% tipo «paleta → tipografia → lockup» salian rotas). Mapeamos el caracter a su
% version matematica, que si tiene glifo. Asi el autor puede seguir escribiendo
% «→» comodamente en el YAML.
\usepackage{newunicodechar}
\newunicodechar{→}{\ensuremath{\rightarrow}}

\setmainfont{Helvetica Neue}
\setsansfont{Helvetica Neue}
\setlength{\parindent}{0pt}
\setlength{\parskip}{6pt}
\hyphenpenalty=200
\exhyphenpenalty=200
\pretolerance=400
\tolerance=2000
\setlength{\emergencystretch}{1em}
\renewcommand{\arraystretch}{1.25}
\setlist{leftmargin=5mm,itemsep=1mm,topsep=1mm}
\newcolumntype{Y}{>{\RaggedRight\arraybackslash}X}

\definecolor{milcMagenta}{HTML}{E6007E}
\definecolor{milcVino}{HTML}{5A0038}
\definecolor{milcTurquesa}{HTML}{0093A5}
\definecolor{milcVerde}{HTML}{58A923}
\definecolor{milcMostaza}{HTML}{E5B400}
\definecolor{milcOcre}{HTML}{B86600}
\definecolor{milcNegro}{HTML}{241816}
\definecolor{milcGris}{HTML}{F7F7F4}
\definecolor{milcLinea}{HTML}{E8E2DD}
\definecolor{milcGradePrimary}{HTML}{FF2D87}
\definecolor{milcGradeDark}{HTML}{9F1B56}
\definecolor{milcGradeSoft}{HTML}{FFE0EE}
\definecolor{milcGradeAccent}{HTML}{CFE7FF}
\definecolor{milcGradeText}{HTML}{FFFFFF}
\color{milcNegro}

\pagestyle{fancy}
\fancyhf{}
\lhead{\small Tecnología e Informática · Grado 10}
\rhead{\small Guía 23 · MILC v3}
\cfoot{\small\thepage}
\renewcommand{\headrulewidth}{0pt}

\newtcolorbox{titlebox}[1]{
  enhanced,colback=#1,colframe=#1,arc=7mm,boxrule=0pt,
  left=7mm,right=7mm,top=3mm,bottom=3mm,
  before skip=8pt,after skip=8pt,
  fontupper=\sffamily\bfseries\Large\color{white}
}

\newtcolorbox{softbox}[3]{
  enhanced,colback=#2,colframe=#1,borderline west={4pt}{0pt}{#1},
  arc=2mm,boxrule=.4pt,left=4mm,right=4mm,top=3mm,bottom=3mm,
  before skip=6pt,after skip=8pt,title={#3},coltitle=white,
  colbacktitle=#1,fonttitle=\sffamily\bfseries,fontupper=\RaggedRight,
  attach boxed title to top left={xshift=3mm,yshift=-2mm},
  boxed title style={arc=2mm, boxrule=0pt}
}

\newtcolorbox{drawbox}[2]{
  enhanced,colback=white,colframe=#1,arc=4mm,boxrule=1pt,
  left=5mm,right=5mm,top=4mm,bottom=4mm,before skip=8pt,after skip=8pt,
  title={#2},coltitle=white,colbacktitle=#1,fonttitle=\sffamily\bfseries,
  fontupper=\RaggedRight,
  attach boxed title to top left={xshift=4mm,yshift=-2mm},
  boxed title style={arc=3mm, boxrule=0pt}
}

\newtcolorbox{infoband}[1]{
  enhanced,colback=#1!8,colframe=#1!50,arc=2mm,boxrule=.5pt,
  left=4mm,right=4mm,top=2mm,bottom=2mm,before skip=4pt,after skip=4pt,
  fontupper=\small\RaggedRight
}

% Puente narrativo entre secciones (contrato editorial Conéctate).
% Da continuidad: "Acabas de X. Ahora vas a Y porque Z."
% Visualmente: banda izquierda en color del grado, texto en itálica pequeña.
\newtcolorbox{puentebox}{
  enhanced,colback=milcGradeSoft,colframe=milcGradeDark,
  borderline west={3pt}{0pt}{milcGradeDark},
  arc=0mm,boxrule=0pt,left=5mm,right=4mm,top=2.5mm,bottom=2.5mm,
  before skip=4pt,after skip=8pt,
  fontupper=\itshape\small\color{milcNegro}\RaggedRight
}
% La flecha va en modo matematico a proposito: Helvetica Neue NO tiene el glifo
% U+2192, asi que \textrightarrow se compilaba a NADA («Missing character: There
% is no → in font Helvetica Neue») y el puente salia sin flecha. En modo
% matematico la flecha viene de la fuente matematica y si se dibuja.
\newcommand{\puente}[1]{%
  \begin{puentebox}%
    \textcolor{milcGradeDark}{$\rightarrow$}\hspace{2mm}#1%
  \end{puentebox}%
}

\newcommand{\linead}[1]{\noindent\color{milcLinea}\rule{#1}{0.5pt}\color{milcNegro}}
\newcommand{\respuesta}[1]{\par\vspace{2mm}\foreach \n in {1,...,#1}{\noindent\color{milcLinea}\rule{\linewidth}{0.5pt}\par\vspace{5mm}}\color{milcNegro}}
\newcommand{\checkbox}{\textcolor{milcGradeDark}{\fbox{\phantom{x}}}\hspace{5pt}}

\newcommand{\phasecard}[4]{
\begin{tcolorbox}[enhanced,colback=#3,colframe=#2,arc=3mm,boxrule=.5pt,height=3.05cm,valign=center,halign=center,left=1.5mm,right=1.5mm,top=2mm,bottom=2mm]
{\sffamily\bfseries\fontsize{22}{24}\selectfont\textcolor{#2}{#1}}\par
{\sffamily\bfseries\small #4}
\end{tcolorbox}
}

% ── Recursos visuales de la guía ──────────────────────────────────────
% Declarados en el bloque `recursos:` del YAML y emitidos por el builder.
% \guiaFigura   → imagen rasterizada o PDF (foto, carta celeste, montaje).
% \guiaDiagrama → fuente TikZ/circuitikz (.tex) incrustada como vector:
%                 es la vía para circuitos y esquemáticos editables.
% #1 = ruta relativa al .tex · #2 = pie (puede ir vacío)
\newcommand{\guiaFigura}[2]{%
\begin{center}
\includegraphics[width=0.88\linewidth,height=0.38\textheight,keepaspectratio]{#1}\\[1.5mm]
{\footnotesize\itshape\textcolor{milcNegro!75}{#2}}
\end{center}
}

\newcommand{\guiaDiagrama}[2]{%
\begin{center}
\input{#1}\\[1.5mm]
{\footnotesize\itshape\textcolor{milcNegro!75}{#2}}
\end{center}
}

\begin{document}
\thispagestyle{empty}
\begin{tikzpicture}[remember picture,overlay]
  \fill[white] (current page.south west) rectangle (current page.north east);
  \path[fill=milcGradePrimary,rounded corners=16mm]
    ([xshift=.55cm,yshift=-.55cm]current page.north west) rectangle
    ([xshift=-.55cm,yshift=3.65cm]current page.south east);

  \node[anchor=north west,text=milcGradeText] at ([xshift=2.0cm,yshift=-1.85cm]current page.north west)
    {\sffamily\bfseries\fontsize{14}{16}\selectfont GRADO};
  \node[anchor=north west,text=milcGradeText,opacity=.82] at ([xshift=2.0cm,yshift=-2.75cm]current page.north west)
    {\sffamily\bfseries\fontsize{9}{11}\selectfont SERIE GUÍAS · TECNOLOGÍA E INFORMÁTICA};

  \begin{scope}[shift={([xshift=-3.05cm,yshift=-2.55cm]current page.north east)}]
    \fill[white,opacity=.80] (-.62,-.52) rectangle (.62,.52);
    \draw[milcGradeText,line width=1pt] (-.62,-.52) rectangle (.62,.52);
    \fill[milcGradeDark] (-.42,-.38) rectangle (-.22,.06);
    \fill[milcGradeAccent] (-.10,-.38) rectangle (.10,.24);
    \fill[milcGradePrimary!60!black] (.22,-.38) rectangle (.42,.38);
  \end{scope}

  \node[anchor=west,text=milcGradeText] at ([xshift=1.9cm,yshift=-12.85cm]current page.north west)
    {\sffamily\bfseries\fontsize{124}{124}\selectfont 10};
  % El circulito de grado (°) solo aplica a las guias de grado («11°»); en
  % Bebras y Semillero el numero grande es un momento/modulo, no un grado.
  \draw[milcGradeText,line width=1.7pt]
    ([xshift=4.52cm,yshift=-11.68cm]current page.north west) circle (.13cm);

  \node[anchor=west,text=milcGradeText,text width=8.7cm,align=left] at ([xshift=10.35cm,yshift=-11.6cm]current page.north west)
    {
     {\sffamily\bfseries\fontsize{19}{22}\selectfont Google Sheets + IA ---\\contabilidad personal\\o microempresa}\\[4mm]
     {\sffamily\bfseries\fontsize{23}{26}\selectfont Guía 23-10-TIC}\\[2mm]
     {\sffamily\fontsize{13}{16}\selectfont Período 3 · Ofimática con IA: contabilidad, datos y emprendimiento}\\[5mm]
     {\sffamily\bfseries\fontsize{11}{13}\selectfont Institución Educativa Sor María Juliana}\\[1mm]
     {\sffamily\fontsize{10}{12}\selectfont Autor: Dr. Álvaro Cárdenas Orozco}
    };

  \path[fill=milcGradeSoft,rounded corners=7mm]
    ([xshift=1.35cm,yshift=.85cm]current page.south west) rectangle
    ([xshift=-1.35cm,yshift=4.25cm]current page.south east);
  \draw[milcGradeDark,line width=.65pt,rounded corners=7mm]
    ([xshift=1.35cm,yshift=.85cm]current page.south west) rectangle
    ([xshift=-1.35cm,yshift=4.25cm]current page.south east);
  \node[anchor=center,text=milcNegro,text width=17.0cm,align=center] at ([yshift=2.55cm]current page.south)
    {\sffamily\fontsize{10}{12}\selectfont Metodología MILC · Apertura ancestral · Pensamiento computacional · Triángulo Dussel-Estoicismo-Floridi\\
    \textcolor{milcGradeDark}{\bfseries Producto:} Plantilla funcional reutilizable de contabilidad en Google Sheets con (a) hoja de Ingresos con tabla y fórmulas, (b) hoja de Egresos con tabla y fórmulas, (c) hoja de Balance con dashboard automático usando SUMA, PROMEDIO, MAX, MIN, SUMAR.SI, (d) al menos 1 gráfico automático de gastos por categoría (circular o barras), (e) integración de tu cuaderno de la sesión 2 migrado a la plantilla + bitácora con al menos 1 duda técnica resuelta con asistencia de IA documentada};
\end{tikzpicture}
\mbox{}
\newpage

\section*{Apertura: Google Sheets + IA para contabilidad personal o microempresa}

\begin{softbox}{milcGradeDark}{milcGradeSoft}{Saber ancestral}
En las oficinas pequeñas del centro de Cartago, en las contadurías de barrio, hubo durante décadas un personaje que cualquier microempresario conocía: \textbf{el contador del barrio que llevaba la planilla cuadriculada}. No era contador titulado de empresa grande: era el señor con escritorio modesto, lápiz, regla, calculadora y \textbf{planillas cuadriculadas} (hojas grandes pre-impresas con columnas y filas). Cada microempresario del barrio le llevaba sus cuadernos de cuentas mensualmente. El contador desplegaba la planilla, organizaba los datos en columnas (ingresos por concepto, egresos por categoría, totales semanales, totales mensuales), hacía los cálculos con la calculadora, registraba con lápiz, sumaba. Al final del mes, entregaba al microempresario \textbf{el resumen}: balance, comparación con mes anterior, recomendaciones simples (\emph{``este mes subieron los gastos en transporte, revise rutas''}). La planilla cuadriculada era \textbf{la hoja de cálculo ancestral}; el contador era el \textbf{Excel humano} del barrio. Hoy, Google Sheets es esa planilla digital al alcance de cualquiera, y la IA es el aprendiz que ayuda cuando te trabas con una fórmula. Pero la regla del oficio sigue siendo la misma: \textbf{tú firmas las cuentas; la IA solo asiste}.
\end{softbox}

\begin{softbox}{milcTurquesa}{milcGris}{Saber contemporáneo}
\textbf{Google Sheets} es la hoja de cálculo gratuita más usada del mundo, integrada con cuenta Gmail. Sus 5 funciones básicas para contabilidad son: (1) \texttt{=SUMA(rango)}: suma valores. Ej. \texttt{=SUMA(B2:B30)} suma todos los egresos del mes. (2) \texttt{=PROMEDIO(rango)}: media aritmética. Ej. \texttt{=PROMEDIO(B2:B30)} promedio de egresos. (3) \texttt{=MAX(rango)} y \texttt{=MIN(rango)}: valor máximo y mínimo. (4) \texttt{=SUMAR.SI(rango criterio; criterio; rango suma)}: suma condicional. Ej. \texttt{=SUMAR.SI(C2:C30; ``Transporte''; B2:B30)} suma todos los egresos cuya categoría sea Transporte. (5) \texttt{=GEMINI("pregunta")}: integración nativa con IA Gemini desde 2025. Permite preguntar a la IA dentro de una celda. Ej. \texttt{=GEMINI("propone 3 categorías para estos egresos:" \&unir(B2:B20;","))}. La regla profesional: \textbf{las fórmulas dan velocidad; la IA da sugerencias; tú decides}. Los \textbf{gráficos automáticos} se generan con Insertar → Gráfico, eligiendo el tipo (barras para comparar, líneas para tendencia, circular para proporciones). La hoja de Balance integra todo: dashboard con totales, promedios, gráfico de gastos por categoría.
\end{softbox}

\begin{softbox}{milcMostaza}{milcGris}{Pregunta-puente}
\textit{¿Qué sabía el contador del barrio al organizar la planilla cuadriculada con disciplina, que el novato olvida cuando intenta hacer Sheets sin estructura? ¿Y por qué \texttt{=GEMINI()} dentro de Sheets es asistente pero nunca decisor final?}
\end{softbox}

\begin{softbox}{milcVerde}{milcGris}{Saber y saber hacer}
Al terminar podrás: (1) \textbf{identificar} las 5 funciones básicas de Sheets útiles para contabilidad personal; (2) \textbf{aplicar} las fórmulas a tus datos del cuaderno de la sesión 2 ya migrados; (3) \textbf{analizar} el balance con gráficos automáticos para detectar patrones; (4) \textbf{evaluar} cuándo usar \texttt{=GEMINI()} y cuándo confiar solo en fórmulas propias.
\end{softbox}



\small
\begin{tabularx}{\textwidth}{>{\bfseries}p{3.0cm}Y}
\rowcolor{milcGradePrimary!12}Autor & Dr. Álvaro Cárdenas Orozco\\
DBA & Aplicación de ofimática asistida por IA a contabilidad básica y emprendimiento (MEN, T\&I)\\
Referentes & Lean Startup · Phronesis financiera · Ética del dato empresarial\\
Producto final & Plantilla funcional reutilizable de contabilidad en Google Sheets con (a) hoja de Ingresos con tabla y fórmulas, (b) hoja de Egresos con tabla y fórmulas, (c) hoja de Balance con dashboard automático usando SUMA, PROMEDIO, MAX, MIN, SUMAR.SI, (d) al menos 1 gráfico automático de gastos por categoría (circular o barras), (e) integración de tu cuaderno de la sesión 2 migrado a la plantilla + bitácora con al menos 1 duda técnica resuelta con asistencia de IA documentada\\
\end{tabularx}
\normalsize

\puente{Antes de empezar, mira el viaje completo. En la sesión 2 llevaste cuaderno de cuentas durante 2 semanas. Hoy migras al \textbf{formato digital con automatización}. Las próximas sesiones (flujo de caja, visualización, mercado, Canvas, comunicación, proyecto) usarán esta plantilla como base.}

\begin{titlebox}{milcGradeDark}
Ruta MILC: así vamos a aprender
\end{titlebox}
\begin{tabularx}{\textwidth}{>{\centering\arraybackslash}X>{\centering\arraybackslash}X>{\centering\arraybackslash}X>{\centering\arraybackslash}X}
\phasecard{1}{milcVerde}{milcGris}{Escucha\\[-1mm]\small Observo, escucho y pregunto.} &
\phasecard{2}{milcTurquesa}{milcGris}{Sistematización\\[-1mm]\small Pensamiento computacional.} &
\phasecard{3}{milcMagenta}{milcGris}{Praxis\\[-1mm]\small Construyo y aplico.} &
\phasecard{4}{milcVino}{milcGris}{Evaluación\\[-1mm]\small Reflexión liberadora.}
\end{tabularx}


\puente{Antes de teorizar, vas a hacer un \emph{primer Sheet}: abre Google Sheets, crea nuevo, y replica una columna de tus ingresos del cuaderno físico. En la celda debajo de los datos, escribe \texttt{=SUMA(B2:B10)} (ajusta el rango). Mira el resultado. Acabas de hacer tu primera fórmula.}

\begin{titlebox}{milcVerde}
Fase 1 · Escucha
\end{titlebox}

\begin{softbox}{milcVerde}{milcGris}{Escena para escuchar}
\textbf{Actividad 1 · IDENTIFICA --- Mi primera fórmula en Sheets} (15 min · individual con dispositivo). Abre Google Sheets (\texttt{sheets.google.com} o desde Drive). Crea hoja nueva. En la columna B, escribe 5 valores de ingresos reales de tu cuaderno (uno por celda, B2 a B6). En la celda B8 escribe la fórmula: \texttt{=SUMA(B2:B6)} y presiona Enter. Mira el resultado: aparece el total de los 5 valores. Después prueba \texttt{=PROMEDIO(B2:B6)} en B9. Acabas de hacer 2 fórmulas básicas.
\end{softbox}

\checkbox Abrí Google Sheets y creé hoja nueva.\\
\checkbox Escribí 5 valores reales en columna.\\
\checkbox Apliqué =SUMA y =PROMEDIO y vi resultados.

\begin{infoband}{milcVerde}


\textbf{Lo que va al cuaderno (Actividad 1 · IDENTIFICA).} Título: «Actividad 1 --- Primera fórmula». \textbf{Formato:} captura de la hoja con las 2 fórmulas funcionando. \textbf{Sabes que terminaste cuando} ves que el cálculo automático es rápido.
\end{infoband}

\puente{Ya tienes primer cálculo. Ahora vas a entender la \textbf{anatomía de Sheets para contabilidad}: 5 fórmulas básicas, estructura de 3 hojas, integración con IA, los 5 errores típicos.}

\begin{titlebox}{milcTurquesa}
Fase 2 · Sistematización con pensamiento computacional
\end{titlebox}

\begin{softbox}{milcTurquesa}{milcGris}{Cuatro pilares aplicados al tema}
Tu primera fórmula te muestra la velocidad de Sheets. Ahora necesitas la \textbf{anatomía profesional} para construir plantilla de contabilidad completa.
\end{softbox}

\small
\begin{tabularx}{\textwidth}{>{\bfseries\RaggedRight\arraybackslash}p{3.6cm}Y}
\rowcolor{milcTurquesa!90}\color{white}Pilar & \color{white}\bfseries Aplicación al tema\\
1. Descomposición & Las \textbf{5 fórmulas básicas} de contabilidad se descomponen así: (1) \texttt{=SUMA(B2:B30)}: suma todos los valores del rango. (2) \texttt{=PROMEDIO(B2:B30)}: media aritmética. (3) \texttt{=MAX(B2:B30)}: valor más alto (mayor gasto). (4) \texttt{=MIN(B2:B30)}: valor más bajo (menor gasto). (5) \texttt{=SUMAR.SI(C2:C30; "Transporte"; B2:B30)}: suma solo cuando la columna C contiene \emph{Transporte}. Esa última es crucial para totalizar por categoría.\\
2. Reconocimiento de patrones & La \textbf{estructura de 3 hojas} para contabilidad personal: (a) \textbf{Hoja Ingresos}: columnas Fecha, Fuente, Monto. (b) \textbf{Hoja Egresos}: columnas Fecha, Categoría, Descripción, Monto. (c) \textbf{Hoja Balance}: dashboard que usa fórmulas \texttt{=SUMAR.SI()} para totalizar por categoría, \texttt{=SUMA()} para totales generales, gráfico automático. Esta estructura es replicable: el mismo formato sirve para contabilidad personal o microempresa.\\
3. Abstracción & La \textbf{integración con IA} en Sheets tiene 2 vías: (a) \texttt{=GEMINI("pregunta")} directamente en celda. Útil para clasificación automática (\emph{=GEMINI("clasifica este gasto: " \&B2)}). Disponible desde 2025 con Workspace. (b) Pegar datos en ChatGPT/Gemini externos pidiendo análisis. Útil para análisis cualitativos. La regla profesional: \textbf{la IA propone, tú verificas con casos conocidos antes de confiar la fórmula}.\\
4. Algoritmo conceptual & Algoritmo: (1) crear nueva hoja de Google Sheets con 3 pestañas. (2) configurar columnas en Ingresos y Egresos. (3) migrar datos del cuaderno físico de la sesión 2. (4) en Balance, aplicar \texttt{=SUMAR.SI()} para totalizar por categoría. (5) calcular balance neto. (6) generar gráfico automático de gastos por categoría. (7) cuando te trabes con alguna fórmula, preguntar a Gemini o ChatGPT. (8) documentar la duda y la respuesta en bitácora.\\
\end{tabularx}
\normalsize

\begin{drawbox}{milcTurquesa}{Anatomía de Sheets contable}
\textbf{5 fórmulas:} SUMA / PROMEDIO / MAX / MIN / SUMAR.SI. \\[1mm] \textbf{3 hojas:} Ingresos / Egresos / Balance. \\[1mm] \textbf{Gráfico:} Insertar → Gráfico automático. \\[1mm] \textbf{IA:} GEMINI integrada o ChatGPT externo. \\[1mm] \textbf{Regla:} IA propone, tú verificas.
\end{drawbox}

\begin{softbox}{milcMostaza}{milcGris}{Errores comunes y su corrección}
(1) \textbf{Rango mal seleccionado}: \texttt{=SUMA(B:B)} cuando solo querías B2:B30 (incluye el encabezado). (2) \textbf{Fórmulas a mano celda por celda}: en lugar de arrastrar la fórmula para que se replique. (3) \textbf{Categorías inconsistentes}: \emph{``Transporte''} y \emph{``transporte''} con minúscula son distintas para SUMAR.SI. (4) \textbf{Aceptar fórmula de IA sin probar}: =GEMINI puede sugerir fórmulas mal estructuradas. (5) \textbf{Sin gráfico}: el dashboard sin visualización pierde su valor.
\end{softbox}

\begin{infoband}{milcTurquesa}


\textbf{Lo que va al cuaderno (Actividad 2 · APLICA).} Título: «Actividad 2 --- Anatomía Sheets contable». \textbf{Formato:} (a) las 5 fórmulas con ejemplo; (b) estructura 3 hojas; (c) los 5 errores con marca. \textbf{Sabes que terminaste cuando} puedes construir plantilla básica sin abrir esta guía.
\end{infoband}

\puente{Tienes el método. Toca aplicarlo: armas la plantilla de 3 hojas (Ingresos, Egresos, Balance), migras los datos del cuaderno, aplicas fórmulas, generas gráfico, integras IA para 1 duda técnica.}

\begin{titlebox}{milcMagenta}
Fase 3 · Praxis con pensamiento computacional
\end{titlebox}

\begin{softbox}{milcMagenta}{milcGris}{Construyendo el producto con los cuatro pilares}
\textbf{Actividad 3 · ANALIZA + EVALÚA --- Plantilla Sheets contable + IA documentada} (40 min · individual con dispositivo). Hora de aplicar. Construyes plantilla de 3 hojas, migras datos, generas gráfico, documentas 1 duda resuelta con IA.
\end{softbox}

\small
\begin{tabularx}{\textwidth}{>{\bfseries\RaggedRight\arraybackslash}p{3.6cm}Y}
\rowcolor{milcMagenta!90}\color{white}Pilar & \color{white}\bfseries Aplicación al producto\\
Descomposición & Tu plantilla se construye en 7 pasos: (1) crear nueva hoja en Sheets. (2) renombrar 3 pestañas: Ingresos, Egresos, Balance. (3) configurar encabezados en Ingresos (Fecha, Fuente, Monto) y Egresos (Fecha, Categoría, Descripción, Monto). (4) migrar los datos de tu cuaderno físico de la sesión 2 (las 2 semanas reales). (5) en la hoja Balance, crear celdas con fórmulas: total ingresos, total egresos, balance, sumas por categoría con \texttt{=SUMAR.SI()}. (6) seleccionar la tabla de sumas por categoría e Insertar → Gráfico (Sheets sugiere circular o barras automáticamente). (7) cuando te trabes con alguna fórmula, preguntar a IA y documentar.\\
Patrones & Sigue el patrón \textbf{``plantilla reutilizable, no solo para hoy''}: tu plantilla debe servir para que el próximo mes solo cambies los datos. Las fórmulas se actualizan solas. Las categorías se mantienen iguales. El gráfico se regenera. Esa replicabilidad es lo que diferencia plantilla profesional de cálculo único.\\
Abstracción & La \textbf{bitácora de IA} documenta al menos 1 duda técnica que tuviste y cómo la IA te ayudó: (a) duda específica (\emph{``¿cómo sumo solo los egresos de transporte sin contar el resto?''}). (b) pregunta hecha a IA. (c) respuesta de IA (la fórmula que sugirió). (d) verificación (probaste con datos conocidos y funcionó o no). (e) decisión final.\\
Algoritmo de armado & Algoritmo: (1) crear plantilla con 3 pestañas. (2) migrar datos. (3) aplicar fórmulas. (4) generar gráfico. (5) bitácora IA. (6) guardar y compartir con tu correo institucional como respaldo.\\
\end{tabularx}
\normalsize

\begin{drawbox}{milcMagenta}{Checklist antes de entregar}
\checkbox Plantilla Google Sheets con 3 pestañas. \\[1mm] \checkbox Datos migrados del cuaderno físico. \\[1mm] \checkbox Fórmulas SUMA, PROMEDIO, MAX, MIN aplicadas. \\[1mm] \checkbox SUMAR.SI por cada categoría de egresos. \\[1mm] \checkbox Balance neto calculado automáticamente. \\[1mm] \checkbox Gráfico automático de gastos por categoría. \\[1mm] \checkbox Bitácora con 1+ duda resuelta con IA documentada.
\end{drawbox}

\begin{softbox}{milcMagenta}{milcGris}{Plantilla de trabajo}
\textbf{Hoja Ingresos.} \textit{Fecha + Fuente + Monto.} \\[1mm] \textbf{Hoja Egresos.} \textit{Fecha + Categoría + Descripción + Monto.} \\[1mm] \textbf{Hoja Balance.} \textit{Sumas + Promedios + MAX/MIN + SUMAR.SI + Balance + Gráfico.} \\[1mm] \textbf{Bitácora IA.} \textit{Duda + Pregunta + Respuesta + Verificación.}
\end{softbox}

\begin{infoband}{milcMagenta}


\textbf{Lo que va al cuaderno (Actividad 3 · ANALIZA + EVALÚA).} Título: «Actividad 3 --- Mi plantilla Sheets». \textbf{Formato:} (a) link al archivo Sheets; (b) capturas de las 3 hojas; (c) bitácora IA. \textbf{Sabes que terminaste cuando} tu plantilla podría servirle a un microempresario.
\end{infoband}

\puente{Lo que sigue resume \emph{qué} entregas hoy: plantilla funcional + datos migrados + gráfico + bitácora IA. El criterio principal: que la plantilla puedas reusarla para los próximos meses sin tener que rehacerla.}

\begin{titlebox}{milcMostaza}
Producto de la sesión
\end{titlebox}
\textbf{Plantilla Sheets 3 hojas + datos + gráfico + bitácora IA}.
\begin{softbox}{milcMostaza}{milcGris}{Criterios de calidad}
(1) Plantilla con 3 pestañas funcionales. (2) Datos reales migrados. (3) 5 fórmulas básicas aplicadas. (4) Gráfico automático generado. (5) Bitácora IA con duda resuelta. (6) Plantilla replicable para meses futuros.
\end{softbox}

\puente{Antes de cerrar, mira Sheets desde las cinco dimensiones humanas. Automatizar cálculos con disciplina es respeto por tu propio tiempo futuro; hacer las cuentas a mano cuando hay herramienta gratuita es nostalgia mal entendida.}

\begin{titlebox}{milcVino}
Fase 4 · Evaluación liberadora — cinco dimensiones
\end{titlebox}

\begin{softbox}{milcVino}{milcGris}{Sentido de la evaluación}
Evaluar no es solo revisar una nota. Es reconocer qué aprendí, qué cambió en mí, qué emociones manejé, cómo me relaciono con la ciudad y la comunidad, y qué le dejaré a quien viene detrás.
\end{softbox}

\textbf{Mi reflexión liberadora en cinco dimensiones.} Completa cada frase iniciada:

\small
\begin{tabularx}{\textwidth}{>{\bfseries\RaggedRight\arraybackslash}p{3.4cm}Y>{\RaggedRight\arraybackslash}p{4.6cm}}
\rowcolor{milcVino}\color{white}Dimensión & \color{white}\bfseries Mi respuesta (frame starter) & \color{white}\bfseries Algo que descubrí\\
1. Personal & Lo que más cambió en mí fue \linead{6.5cm} & \linead{4.2cm}\\
2. Emocional & La emoción más fuerte fue \linead{2.5cm} y la regulé \linead{3cm} & \linead{4.2cm}\\
3. Ciudadana & En democracia esto importa porque \linead{6cm} & \linead{4.2cm}\\
4. Local (Cartago) & En mi barrio sería útil para \linead{6.5cm} & \linead{4.2cm}\\
5. Intergeneracional & A mi abuela/abuelo le diría \linead{6.5cm} & \linead{4.2cm}\\
\end{tabularx}
\normalsize

\begin{infoband}{milcVino}
\textbf{Para la próxima sustentación.} Vuelve a esta tabla en seis meses. Verás que las respuestas cambian — no porque lo aprendido cambie, sino porque tú cambias. La evaluación liberadora es una conversación contigo a través del tiempo.
\end{infoband}


\puente{Y para terminar, tres voces sobre la planilla digital. Dussel sobre las herramientas que liberan al microempresario, Marco Aurelio sobre la disciplina de la fórmula correcta, Floridi sobre la IA como copiloto contable.}

\begin{titlebox}{milcGradeDark}
Triángulo de pensamiento · cierre filosófico
\end{titlebox}

\begin{softbox}{milcVino}{milcGris}{Dussel · lente del nosotros}
\textit{``Las herramientas digitales gratuitas con IA democratizan el oficio contable; el microempresario que las domina no depende de contadores caros.''}

Tradicionalmente, los microempresarios pagaban contadores para tareas que ahora pueden hacer con Sheets gratis y asistencia de IA. La filosofía de la liberación pide aprovechar esa democratización: \textbf{la tecnología accesible es libertad económica}. La phronesis del oficio digital consiste en \textbf{usar las herramientas gratuitas con disciplina para no depender de servicios pagos innecesarios}.

\textbf{¿Estoy aprendiendo Sheets como acto de libertad económica futura, o solo como tarea escolar?}
\end{softbox}
\linead{\linewidth}\\[1mm]
\linead{\linewidth}

\begin{softbox}{milcMostaza}{milcGris}{Estoicismo · Marco Aurelio · cuidado interior}
\textit{``La fórmula correcta es virtud del oficio; copiar fórmula sin entenderla es vanidad técnica.''}

Es tentador copiar fórmulas que la IA sugiere sin entenderlas. La sabiduría estoica recomienda lo opuesto: \emph{entiende cada fórmula antes de aceptarla}. Solo entendiendo qué hace =SUMAR.SI puedes detectar si la IA la sugirió mal. La phronesis del oficio consiste en \textbf{usar la IA para aprender, no para reemplazar el aprendizaje}.

\textbf{¿Entiendo cada fórmula de mi plantilla, o solo copié lo que la IA dijo?}
\end{softbox}
\linead{\linewidth}\\[1mm]
\linead{\linewidth}

\begin{softbox}{milcTurquesa}{milcGris}{Floridi · lente de la infoesfera}
\textit{``La IA como copiloto contable es ética profesional en la era de la automatización accesible.''}

Tu plantilla con \texttt{=GEMINI()} integrada es modelo del oficio contemporáneo: herramienta gratuita + IA como asistente + criterio humano para verificar. Esa combinación es ética porque produce resultados confiables sin depender de empresas caras. La phronesis profesional del siglo XXI consiste en \textbf{combinar herramientas accesibles con criterio propio}.

\textbf{¿Mi plantilla combina herramienta gratuita con criterio propio, o pretendo que solo la IA decida?}
\end{softbox}
\linead{\linewidth}\\[1mm]
\linead{\linewidth}

\begin{titlebox}{milcVino}
Compromiso final
\end{titlebox}

\small
\begin{tabularx}{\textwidth}{>{\RaggedRight\arraybackslash}p{3.0cm}YYY}
\rowcolor{milcVino}\color{white}\bfseries Criterio & \color{white}\bfseries Lo logré & \color{white}\bfseries En proceso & \color{white}\bfseries Necesito apoyo\\
Comprensión del tema & Explico ideas centrales con mis palabras. & Reconozco ideas pero falta precisión. & Necesito repasar conceptos base.\\
Pensamiento computacional & Apliqué los 4 pilares al diseñar mi producto. & Apliqué algunos pilares. & Necesito releer la fase 2.\\
Aplicación y producto & Mi producto cumple los criterios y se puede revisar. & Producto funcional con ajustes pendientes. & Aún en construcción.\\
Reflexión 5 dimensiones & Respondí cada dimensión con honestidad. & Respondí algunas con detalle. & Mis respuestas son muy generales.\\
Triángulo de pensamiento & Conecté las 3 voces con mi trabajo real. & Conecté al menos una. & No relaciono las voces con mi proceso.\\
\end{tabularx}
\normalsize

\textbf{Hoy entendí que:}\\[1mm]
\linead{\linewidth}\\[1mm]
\linead{\linewidth}

\textbf{Mi compromiso para la próxima sesión es:}\\[1mm]
\linead{\linewidth}\\[1mm]
\linead{\linewidth}

\begin{softbox}{milcMostaza}{milcGris}{Cita de despedida · Marco Aurelio}
\textit{``El obstáculo en el camino se vuelve el camino. Lo que estorba al camino se vuelve camino.''}\\[1mm]
Cada error de hoy, bien observado, se convierte en pieza del camino que pisarás mañana.
\end{softbox}

\small
\begin{tabularx}{\textwidth}{>{\bfseries}p{3.6cm}Y}
\rowcolor{milcGradeDark!12}Nombre completo: & \linead{10cm}\\
Fecha: & \linead{4cm} \quad Curso/grupo: \linead{4cm}\\
Mi firma: & \linead{9cm}\\
\end{tabularx}
\normalsize

\begin{infoband}{milcGradeDark}
\textbf{Para la próxima sesión.} Llega con tu plantilla Sheets funcional con los datos migrados y el gráfico generado. En la sesión 4 vas a aprender \textbf{flujo de caja y punto de equilibrio}: proyectar 3 meses adelante con la plantilla y calcular cuándo un microemprendimiento se equilibra. Los datos de hoy son la base; la proyección de mañana es la herramienta de decisión.
\end{infoband}

\end{document}
